\documentclass[conference]{IEEEtran}

% ── Packages ──────────────────────────────────────────────────────────
\usepackage[T1]{fontenc}
\usepackage[utf8]{inputenc}
\usepackage{graphicx}
\usepackage{booktabs}
\usepackage{amsmath}
\usepackage{cite}
\usepackage{hyperref}
\usepackage{xcolor}
\usepackage{tabularx}
\usepackage{array}
\usepackage{textcomp}
\usepackage{enumitem}
\usepackage{multirow}
\usepackage{algorithm}
\usepackage{algorithmic}

\hypersetup{
  colorlinks=true,
  linkcolor=black,
  urlcolor=blue,
  citecolor=black
}

% ══════════════════════════════════════════════════════════════════════
\begin{document}

\title{Agentic Finance System: An AI-Powered Personal Finance and Market Insight Platform with Multi-Agent Architecture}

\author{
\IEEEauthorblockN{Jainithissh S}
\IEEEauthorblockA{CB.SC.U4AIE23129}
\and
\IEEEauthorblockN{Krishna Prakash S}
\IEEEauthorblockA{CB.SC.U4AIE232138}
\and
\IEEEauthorblockN{Nitheshkummar C}
\IEEEauthorblockA{CB.SC.U4AIE23155}
\and
\IEEEauthorblockN{Akhilesh Kumar S}
\IEEEauthorblockA{CB.SC.U4AIE232170}
\and
\IEEEauthorblockN{Adarsh Pradeep}
\IEEEauthorblockA{CB.SC.U4AIE23109}
}

\maketitle

% ── Abstract ──────────────────────────────────────────────────────────
\begin{abstract}
We present the Agentic Finance System, an end-to-end AI-powered platform that unifies personal transaction data, real-time market sentiment, and curated financial domain knowledge into explainable, personalised financial guidance. The system employs a multi-agent architecture with three autonomous agents operating on a novel Two-Tier Pipeline: Tier~1 provides deterministic ML/rule-based inference, while Tier~2 layers LLM-powered agentic reasoning via Groq LLaMA~3.3 70B with automatic fallback. The platform integrates FinBERT-based NLP sentiment analysis, Random Forest risk and trend classifiers, a 32-chunk TF-IDF RAG engine, and a Bayesian Knowledge Tracing micro-learning system. An Android client built with Kotlin and Jetpack Compose captures bank SMS automatically via NotificationListenerService. The backend exposes 17 REST endpoints through FastAPI, with graceful degradation ensuring 100\% core feature availability regardless of external API status. The system demonstrates that multi-agent orchestration with deterministic fallbacks can deliver production-grade financial intelligence accessible to students and early-career earners.
\end{abstract}

\begin{IEEEkeywords}
Multi-Agent Systems, Personal Finance, Natural Language Processing, FinBERT, Retrieval-Augmented Generation, Bayesian Knowledge Tracing, Large Language Models
\end{IEEEkeywords}

% ══════════════════════════════════════════════════════════════════════
%  I — INTRODUCTION & PROBLEM STATEMENT
% ══════════════════════════════════════════════════════════════════════
\section{Introduction}

India's digital payments ecosystem processes billions of transactions annually through UPI, generating a massive volume of bank SMS and push notifications. Despite this wealth of personal financial data flowing through users' devices daily, most individuals lack automated tools to convert these messages into actionable spending insights. Financial news arrives as a separate, disconnected stream with no link to an individual's actual financial situation. The result is a fragmented information landscape where personal finance data, market intelligence, and financial education exist in silos.

The emergence of Large Language Models (LLMs) and transformer-based NLP has opened new possibilities for financial AI. However, deploying LLMs in financial applications presents unique challenges: hallucination risks in monetary advice, latency constraints for real-time interaction, and the need for deterministic fallbacks when external APIs become unavailable. A purely LLM-driven approach is insufficient---financial applications demand reliability guarantees that probabilistic models alone cannot provide.

\subsection{Problem Statement}

Existing solutions address fragments of the personal finance problem but none provide a unified, intelligent platform:

\begin{itemize}[leftmargin=*]
  \item \textbf{Expense Trackers} (Walnut, Money Manager): Rely on manual data entry or basic SMS parsing. They store transactions as static records and offer no analytical intelligence, risk assessment, or forward-looking guidance.
  \item \textbf{Market Applications} (Moneycontrol, ET Markets): Deliver generic financial news and market data without any personalisation to the user's actual income, spending, or risk profile.
  \item \textbf{Robo-Advisors} (Groww, Kuvera): Focus narrowly on mutual-fund selection and portfolio management but ignore daily spending patterns, savings behaviour, and financial literacy.
  \item \textbf{Professional Financial Advisors}: Remain inaccessible for students and early-career earners due to high consultation fees and minimum investment thresholds, creating a financial guidance gap for the demographic that needs it most.
  \item \textbf{Financial Literacy Platforms} (Zerodha Varsity, ET Learn): Offer static, one-size-fits-all content without adaptation to the learner's knowledge state or real financial situation.
\end{itemize}

\textbf{The core gap:} No single system unifies \emph{personal transaction data}, \emph{real-time market sentiment}, and \emph{curated financial knowledge} into agentic, explainable, personalised financial guidance with production-grade reliability. Furthermore, no existing platform connects financial education to the user's actual financial profile, ensuring learning relevance.

\subsection{Proposed Solution}

The \textbf{Agentic Finance System} addresses this gap through five integrated capabilities:
\begin{enumerate}[leftmargin=*]
  \item Automatic ingestion and structuring of personal finance data from bank SMS and UPI notifications via a backend-authoritative parsing pipeline.
  \item ML-driven personal risk and spending profile computation using Random Forest classifiers with month-proration for temporal robustness.
  \item Real-time financial news understanding and market trend prediction using FinBERT-based NLP with sector-level entity extraction.
  \item Grounded, explainable recommendation generation via a RAG-augmented multi-agent LLM pipeline with deterministic fallbacks.
  \item Adaptive financial literacy education through a Bayesian Knowledge Tracing micro-learning engine with profile-aware curriculum compilation and gamification.
\end{enumerate}

% ══════════════════════════════════════════════════════════════════════
%  II — RELATED WORK
% ══════════════════════════════════════════════════════════════════════
\section{Related Work}

\subsection{Personal Finance Management Tools}

Traditional expense-tracking applications follow a data-entry paradigm where users manually log transactions. While some apps parse bank SMS, they typically stop at categorisation and offer no analytical intelligence. These tools treat financial data as static records rather than inputs to an analytical pipeline. None integrate real-time market context or provide LLM-powered conversational interfaces grounded in the user's actual financial data.

\subsection{Sentiment Analysis in Finance}

Financial sentiment analysis has evolved from lexicon-based approaches to transformer models. FinBERT, a BERT variant fine-tuned on financial text, achieves state-of-the-art performance on financial sentiment benchmarks. However, most deployments use FinBERT in isolation for market analysis without connecting sentiment signals to individual user portfolios or spending patterns. Our system bridges this gap by feeding Agent~B's sentiment outputs directly into Agent~C's personalised recommendation pipeline.

\subsection{Multi-Agent Systems}

Multi-agent architectures have gained traction in AI systems where specialised agents handle distinct subtasks. In financial AI, agents have been used for portfolio optimisation, trading strategy execution, and risk assessment. Our contribution is the Two-Tier Pipeline design, where each agent maintains both a fast deterministic tier and an LLM-powered reasoning tier, ensuring graceful degradation---a pattern not commonly seen in existing multi-agent financial systems.

\subsection{Adaptive Learning Systems}

Bayesian Knowledge Tracing (BKT), originally developed for intelligent tutoring systems, models learner knowledge as a latent variable updated through observed performance. While BKT has been extensively applied in mathematics and programming education, its application to personal finance literacy is novel. Our system extends classical BKT with a 3-state belief model (Unknown, Partial, Mastered) and latency-aware updates that distinguish between confident knowledge and lucky guesses.

% ══════════════════════════════════════════════════════════════════════
%  III — SYSTEM ARCHITECTURE
% ══════════════════════════════════════════════════════════════════════
\section{System Architecture}

The system follows a \textbf{multi-agent architecture} with three autonomous agents coordinated through a FastAPI backend. Each agent operates on a \textbf{Two-Tier Pipeline}: Tier~1 provides fast, deterministic ML/rule-based inference; Tier~2 adds LLM-powered agentic reasoning via Groq (LLaMA~3.3 70B), with automatic fallback to Tier~1 if the LLM is unavailable.

\subsection{High-Level Layers}

The architecture comprises six layers:
\begin{itemize}[leftmargin=*]
  \item \textbf{Client Layer:} Android app (Kotlin, Jetpack Compose) with background \texttt{NotificationListenerService} for automatic SMS capture. 30~Kotlin source files across 6~screens: Login, Dashboard, Chat, Learn, Roadmap, and Market. Kotlin Coroutines ensure non-blocking HTTP calls via Retrofit~2.
  \item \textbf{API Gateway:} FastAPI backend (Python~3.11, Uvicorn) with Firebase JWT authentication. 17~REST endpoints organised into 6~modular routers with CORS middleware for cross-origin access.
  \item \textbf{Agent Layer:} Agent~A (Finance Analyzer), Agent~B (Market Intelligence), Agent~C (Decision Synthesizer). Each agent is self-contained with dedicated ML models and fallback logic.
  \item \textbf{Learning Layer:} Bayesian Knowledge Tracing engine with AI-generated micro-learning cards, gamified roadmap, and achievement system.
  \item \textbf{Persistence Layer:} Dual-database strategy with Supabase PostgreSQL (production) and SQLite (development), configured via environment variable routing.
  \item \textbf{External APIs:} Groq Cloud LLM (30s timeout), Economic Times RSS feed, NSE live market data (pre-warmed at startup).
\end{itemize}

\subsection{Agent A --- Personal Finance Analyzer}

Agent~A processes the complete lifecycle from raw SMS ingestion to personalised financial risk assessment. The SMS parsing pipeline employs a 3-tier regex extraction cascade:

\textbf{Tier~1 --- Currency Regex:} Matches patterns like ``Rs'', ``INR'', or the rupee symbol followed by a numeric amount with optional comma separators. This handles the most common Indian bank SMS formats.

\textbf{Tier~2 --- Keyword Proximity:} When no currency prefix is found, the parser searches for amounts within 15~characters of transactional keywords such as ``debited'', ``credited'', ``spent'', or ``received''.

\textbf{Tier~3 --- Digit Fallback:} As a last resort, any standalone digit sequence with two or more characters is extracted. Credit/debit classification uses 15~credit keywords versus 12~debit keywords.

Category assignment combines a keyword map (40+ merchant-to-category mappings across 10~classes) with a TF-IDF + Logistic Regression model at confidence threshold $\geq 0.50$. Financial message validation requires the presence of a monetary amount plus either a debit/credit keyword or a bank signal term (``a/c'', ``account'', ``upi'', ``neft'', ``atm'', ``balance''). Messages with $\geq 2$~spam keywords and no bank signals are rejected.

Risk classification employs a Random Forest model (120~trees, depth~6, balanced subsample weighting) trained on 900~synthetic profiles. The model operates on a 6-feature vector: monthly income, fixed expenses, variable expenses, savings, savings rate, and variable share. A month-proration step projects partial-month data to full-month estimates using $\text{projected} = \text{actual} / (\text{days\_covered} / 30)$, preventing early-month low-risk bias. The output includes three confidence bands: high ($\geq 0.75$), medium (0.50--0.75), and low ($< 0.50$).

The LLM tier generates a 4--5 sentence personalised narrative using Groq LLaMA~3.3 70B with a system prompt that enforces rupee-denominated, data-specific answers referencing the user's actual amounts, top spending categories, and savings rate.

\subsection{Agent B --- News \& Market Intelligence}

Agent~B performs three-stage market analysis: headline acquisition, sentiment classification, and trend prediction.

\textbf{Headline Acquisition:} Live Economic Times RSS feed is fetched and parsed (XML), with automatic fallback to a 10-article demo feed if the network is unavailable. Thirteen boilerplate-removal regex patterns strip editorial noise, advertisements, and formatting artifacts before NLP inference.

\textbf{Sentiment Classification:} Each headline is scored using FinBERT (ProsusAI/finbert), a BERT model fine-tuned on 4,500+ financial documents. The model outputs probabilities across three labels: positive, negative, and neutral. A weighted scoring scheme applies 40\% weight to headlines and 60\% to article descriptions for composite sentiment. When FinBERT is unavailable, a keyword-based fallback uses 40+ positive keywords (rally, growth, record high, bullish, surge), 30+ negative keywords (crash, loss, bearish, plunge, decline), and strong-phrase detection with 3$\times$ priority multiplier (``hits record high'', ``crashes'', ``misses estimates'').

\textbf{Named Entity Recognition:} spaCy (en\_core\_web\_sm) extracts company entities, mapped to seven sectors (Banking, IT, Auto, Telecom, FMCG, Metals, Pharma) via a 25+ entry lookup covering companies like Reliance (Energy), HDFC/ICICI (Banking), and TCS/Infosys (IT).

\textbf{Trend Prediction:} A Random Forest model (120~trees, depth~5) aggregates per-article sentiments into a 4-dimensional feature vector: positive\_ratio, negative\_ratio, neutral\_ratio, and net\_score ($= \text{positive\_ratio} - \text{negative\_ratio}$). The model predicts market direction as bullish, sideways, or bearish.

\subsection{Agent C --- Decision Synthesizer \& RAG Chatbot}

Agent~C provides a multi-turn conversational interface that retrieves top-3 relevant chunks from a curated 32-chunk financial knowledge corpus. The corpus covers 12~topic areas: emergency funds, budgeting (50/30/20 rule), debt management (avalanche/snowball methods), investing basics (SIPs, diversification), risk management, tax \& insurance (Section~80C), Indian market specifics (UPI, FDs, Nifty), behavioural finance, goal-based planning, scam awareness, salary optimisation, and real estate decisions.

The custom TF-IDF retrieval algorithm tokenises queries via lowercase word-boundary splitting with 131~stop-word removal, computes log-smoothed IDF scores ($\text{IDF} = \log\frac{n+1}{\text{freq}+1} + 1$), and applies a 3$\times$ bonus for exact tag matches. Retrieved context, combined with Agent~A's financial profile, Agent~B's sentiment signals, and full conversation history, is passed to Groq LLaMA~3.3 70B. The system prompt enforces rupee-denominated, data-grounded answers.

A deterministic keyword-matched fallback covers 6~core topics (emergency funds, budgeting, debt, investing, insurance, tax saving) when the LLM is unavailable, ensuring the chatbot always returns useful guidance.

% ══════════════════════════════════════════════════════════════════════
%  IV — TECHNOLOGY STACK
% ══════════════════════════════════════════════════════════════════════
\section{Technology Stack}

Table~\ref{tab:techstack} summarises the technology stack across all layers of the system.

\begin{table}[!t]
\centering
\caption{Technology Stack}
\label{tab:techstack}
\renewcommand{\arraystretch}{1.15}
\footnotesize
\begin{tabular}{p{1.4cm} p{2.8cm} p{3cm}}
\toprule
\textbf{Layer} & \textbf{Technology} & \textbf{Purpose} \\
\midrule
Backend       & FastAPI + Uvicorn (Python~3.11)     & Async REST API \\
Database      & Supabase PostgreSQL / SQLite        & Dual-DB with failover \\
Auth          & Firebase Admin SDK                  & JWT verification \\
ML -- Risk    & scikit-learn Random Forest          & Risk classification \\
ML -- Cat.    & TF-IDF + Logistic Regression        & SMS categorisation \\
ML -- Trend   & Random Forest on FinBERT agg.       & Market direction \\
NLP -- Sent.  & HuggingFace FinBERT                 & Sentiment scoring \\
NLP -- NER    & spaCy en\_core\_web\_sm              & Entity extraction \\
LLM           & Groq LLaMA~3.3 70B                  & Reasoning \& chatbot \\
RAG           & Custom TF-IDF + IDF                 & Grounded responses \\
Mobile        & Kotlin + Jetpack Compose            & Android client \\
Deploy        & Docker + Render.com                 & Cloud deployment \\
\bottomrule
\end{tabular}
\end{table}

\subsection{ML Models}

Table~\ref{tab:models} details the six ML/AI models employed across the three agents and learning engine.

\begin{table}[!t]
\centering
\caption{ML and AI Models Summary}
\label{tab:models}
\renewcommand{\arraystretch}{1.15}
\footnotesize
\begin{tabular}{p{1.2cm} p{2.6cm} p{1.4cm} p{1.8cm}}
\toprule
\textbf{Model} & \textbf{Algorithm} & \textbf{Data} & \textbf{Output} \\
\midrule
Category  & TF-IDF + Log. Reg.            & 1K+ synth. SMS     & 10 categories \\
Risk      & Random Forest (120, d=6)      & 900 profiles       & low/med/high \\
Trend     & Random Forest (120, d=5)      & 600 regimes        & bull./side./bear. \\
FinBERT   & BERT (fin. fine-tuned)        & ProsusAI           & pos./neg./neut. \\
spaCy     & CNN + trans.\ parser           & OntoNotes 5.0      & ORG/PER/LOC \\
LLaMA     & 70B autoregressive            & Web-scale          & NL generation \\
\bottomrule
\end{tabular}
\end{table}

\subsection{Training Data Generation}

All ML models are trained on carefully designed synthetic datasets to avoid privacy concerns with real financial data, while maintaining statistical properties consistent with Indian financial patterns.

The \textbf{risk model} training data comprises 900~profiles (300 per risk class) with realistic income ranges (15,000--80,000 INR). Fixed expenses are clipped to 20--60\% of income, variable expenses range from 10--120\% of income, and derived features (savings rate, variable share) are computed to ensure internal consistency. The class-balanced sampling with stratified 75/25~train/validation split ensures robust generalisation across risk categories.

The \textbf{category model} uses 1,000+ synthetic SMS messages mimicking formats from major Indian banks (SBI, HDFC, ICICI, Axis) with realistic merchant names, amounts, and transaction descriptions across 10~categories: Food \& Dining, Transport, Shopping, Bills \& Utilities, Groceries, Entertainment, Health, Education, Investment, and Other.

The \textbf{trend model} generates 600~market regime samples (200 per class) with sentiment distributions matching historical market conditions: bullish regimes have high positive ratios (0.4--0.9), bearish have high negative ratios (0.4--0.9), and sideways show balanced sentiment with $\pm 0.15$ skew tolerance.

% ══════════════════════════════════════════════════════════════════════
%  V — FEATURES
% ══════════════════════════════════════════════════════════════════════
\section{Features}

\subsection{Automatic SMS Transaction Parsing}

The system intercepts bank SMS and UPI notifications via Android's \texttt{NotificationListenerService}. The Android client sends \emph{all} notification text to the backend without client-side filtering, placing parsing intelligence entirely server-side. This backend-authoritative design enables rapid iteration on parsing rules without requiring app updates, maintains a single source of truth for classification, and ensures consistent behaviour across all Android devices.

The 3-tier extraction pipeline handles the wide diversity of Indian bank message formats. Credit/debit classification uses 15~credit keywords versus 12~debit keywords. Merchant names are extracted via regex pattern matching (e.g., ``at [merchant name]'' from structured SMS). Category assignment combines a keyword map with 40+ merchant-to-category mappings with ML model refinement at confidence threshold $\geq 0.50$ across 10~classes.

\subsection{Personal Financial Risk Analysis}

Agent~A aggregates transactions over a configurable date range and computes a 6-feature vector for risk assessment. Table~\ref{tab:risk} shows representative risk profiles across the three classification levels.

\begin{table}[!t]
\centering
\caption{Risk Classification Examples (Monthly, INR)}
\label{tab:risk}
\renewcommand{\arraystretch}{1.15}
\footnotesize
\begin{tabular}{p{1.3cm} p{0.9cm} p{0.9cm} p{0.9cm} p{0.9cm} p{0.9cm}}
\toprule
\textbf{Risk} & \textbf{Income} & \textbf{Fixed} & \textbf{Variable} & \textbf{Savings} & \textbf{Rate} \\
\midrule
Low    & 50K & 15K & 10K & 25K & 50\% \\
Medium & 50K & 20K & 20K & 10K & 20\% \\
High   & 50K & 25K & 35K & --10K & --20\% \\
\bottomrule
\end{tabular}
\end{table}

The Random Forest classifier outputs risk level with softmax confidence probability and a confidence band indicator (high/medium/low). Month-proration projects partial-month data to full-month estimates, preventing early-month low-risk bias. Rule-based risk explanations provide actionable advice: low-risk users are advised to ``keep an eye on 1--2 categories'', medium-risk users to ``nudge 5--10\% of variable spends to savings'', and high-risk users to ``pause 1--2 non-essential expenses''. The LLM tier enriches these with personalised narratives referencing the user's actual rupee amounts and top spending categories.

Risk assessment results are logged to the \texttt{risk\_logs} database table with timestamps, date ranges, financial summaries, and both heuristic and ML risk levels, enabling longitudinal analysis of a user's financial health trajectory.

\subsection{News \& Market Sentiment Intelligence}

Agent~B provides real-time market intelligence through a pipeline of headline acquisition, FinBERT sentiment analysis, entity extraction, and trend prediction. NSE index data (Nifty~50, Sensex, Bank~Nifty) is pre-warmed at server startup, eliminating cold-start latency for market data requests. The sentiment analysis applies weighted scoring (40\% headline, 60\% description) to account for the richer context in article descriptions versus clickbait-prone headlines.

Sector-level sentiment aggregation groups articles by extracted company entities, enabling sector-specific market insights (e.g., ``Banking sector sentiment is bullish with 70\% positive headlines''). The trend prediction model aggregates individual article sentiments into a market-direction forecast, providing users with a synthesised market outlook rather than raw headline lists.

\subsection{RAG-Grounded AI Chatbot}

Agent~C provides a multi-turn conversational interface grounded in three knowledge sources: the 32-chunk financial knowledge corpus, Agent~A's financial profile, and Agent~B's market sentiment. The knowledge corpus covers 12~topic areas from emergency funds and budgeting to Indian market specifics and scam awareness.

Each knowledge chunk includes a title, 150--300 word explanation, and curated keyword tags. The TF-IDF retrieval algorithm with tag-match bonus ($\times 3$) achieves topical precision that pure vector-similarity approaches often miss for financial domain queries.  If no chunks achieve a non-zero relevance score, the system returns the first $k$ chunks as generic financial guidance. The deterministic fallback covers 6~core topics (emergency funds, budgeting, debt management, investing basics, insurance fundamentals, tax saving strategies) when the LLM is unavailable, ensuring the chatbot always returns useful guidance.

\subsection{Adaptive Micro-Learning System}

The learning engine implements a probabilistic mastery model using an extended Bayesian Knowledge Tracing framework. Each concept's knowledge state is represented as a 3-state belief vector:
\begin{equation}
P(\text{Unknown}) + P(\text{Partial}) + P(\text{Mastered}) = 1.0
\end{equation}

Belief updates are conditioned on both correctness and response latency, enabling the system to distinguish between different learning states. Table~\ref{tab:bkt} summarises the update rules.

\begin{table}[!t]
\centering
\caption{Bayesian Belief Update Rules}
\label{tab:bkt}
\renewcommand{\arraystretch}{1.15}
\footnotesize
\begin{tabular}{p{1.5cm} p{1.2cm} p{1.1cm} p{1.1cm} p{1.3cm}}
\toprule
\textbf{Outcome} & \textbf{Time} & $\Delta$\textbf{Mast.} & $\Delta$\textbf{Part.} & $\Delta$\textbf{Unkn.} \\
\midrule
Correct & $<30$s & +0.30 & +0.10 & --0.40 \\
Correct & $>30$s & +0.15 & +0.15 & --0.30 \\
Wrong   & $<60$s & --0.30 & +0.10 & +0.20 \\
Wrong   & $>60$s & --0.20 & --0.20 & +0.40 \\
\bottomrule
\end{tabular}
\end{table}

The rationale is: a fast correct answer ($<30$s) strongly indicates mastery, while a slow correct answer may indicate partial understanding or careful reasoning. A fast incorrect answer suggests guessing, while a slow incorrect answer ($>60$s) reveals genuine struggle with the concept.

After each update, belief components are clamped to $[0, 1]$ and renormalized to maintain the probability constraint. Mastery level classification uses thresholds: mastered ($P_{\text{mastered}} > 0.6$), partial ($P_{\text{partial}} > 0.5$), and unknown (otherwise).

\subsection{Concept Graph \& Curriculum Compilation}

A DAG-structured concept graph with 11~core concepts and 4~difficulty levels ensures prerequisite-based learning progression. Foundation concepts (Money Basics, Income Basics) have no prerequisites. Intermediate concepts (Budgeting, Expense Tracking) require foundations. Advanced concepts (Investment Basics, Credit Score, Financial Goals) require multiple intermediate prerequisites. DFS-based cycle detection validates the graph at initialisation.

The curriculum compiler scores candidate concepts using a composite function:
\begin{equation}
\text{score} = \text{readiness} \times \text{urgency} \times \text{relevance}
\end{equation}

\textbf{Readiness} is the minimum mastery score across all prerequisites ($\geq 0.6$ threshold). \textbf{Urgency} is $1 - P_{\text{mastered}}$, prioritising concepts the user has not yet learned. \textbf{Relevance} is a context-based multiplier informed by the user's financial profile: budget/emergency concepts receive $1.5\times$ weight for high-spending-risk users, debt management receives $1.6\times$ for users with debt, and saving strategies receive $1.4\times$ for users with savings rate $< 20\%$. This ensures the learning path adapts not only to knowledge state but also to the user's real financial situation.

AI-generated micro-learning cards (150--200 words + MCQ with 4~options) are produced by Groq LLaMA~3.3 70B at temperature 0.7 for content and 0.8 for quizzes. Cards are cached per mastery level and invalidated when the user's mastery classification changes.

\subsection{Gamified Learning Roadmap}

Five structured lessons progress from \textit{Financial Basics} through \textit{Investment Basics} with prerequisite gating. Table~\ref{tab:lessons} details the learning roadmap.

\begin{table}[!t]
\centering
\caption{Learning Roadmap Structure}
\label{tab:lessons}
\renewcommand{\arraystretch}{1.15}
\footnotesize
\begin{tabular}{p{0.4cm} p{2.2cm} p{0.8cm} p{0.8cm} p{1.8cm}}
\toprule
\textbf{\#} & \textbf{Lesson} & \textbf{Time} & \textbf{Cards} & \textbf{Level} \\
\midrule
1 & Financial Basics  & 10m & 5  & Beginner \\
2 & Budgeting Mastery & 15m & 7  & Beginner \\
3 & Emergency Fund    & 12m & 6  & Intermediate \\
4 & Debt Management   & 20m & 8  & Intermediate \\
5 & Investment Basics & 25m & 10 & Advanced \\
\bottomrule
\end{tabular}
\end{table}

Points are awarded as: 50~base points per lesson, +100~for perfect score, +20~for speed completion ($<5$~minutes). Star ratings follow: 3-star ($\geq 90\%$), 2-star ($\geq 70\%$), 1-star ($\geq 50\%$). Six achievements (First Steps, Bookworm, Perfect Score, Speed Learner, Week Warrior, Master) and daily goals (1~lesson, 10~questions, 15~minutes) with a weighted progress bar ($0.4 + 0.3 + 0.3$) drive engagement.

% ══════════════════════════════════════════════════════════════════════
%  VI — API DESIGN & IMPLEMENTATION
% ══════════════════════════════════════════════════════════════════════
\section{API Design \& Implementation}

The backend exposes 17~REST endpoints, all requiring Firebase Bearer JWT authentication (development mode accepts unauthenticated requests as \texttt{default\_user}). Table~\ref{tab:api} lists all endpoints grouped by agent.

\begin{table}[!t]
\centering
\caption{API Endpoint Reference (17 Endpoints)}
\label{tab:api}
\renewcommand{\arraystretch}{1.1}
\scriptsize
\begin{tabular}{l l l p{2.4cm}}
\toprule
\textbf{Method} & \textbf{Endpoint} & \textbf{Agent} & \textbf{Description} \\
\midrule
POST   & /api/parse\_message        & --    & Parse SMS to txn \\
GET    & /api/transactions          & --    & List transactions \\
DELETE & /api/transactions          & --    & Delete all txns \\
PUT    & /api/trans./\{id\}/cat.    & --    & Override category \\
POST   & /api/analyze\_finance      & A     & Risk analysis + LLM \\
POST   & /api/expense\_breakdown    & A     & Category breakdown \\
GET    & /api/risk\_logs\_summary    & A     & Risk history \\
POST   & /api/analyze\_news         & B     & Sentiment analysis \\
GET    & /api/news\_feed            & B     & Live RSS + FinBERT \\
GET    & /api/live\_market          & B     & NSE indices \\
POST   & /api/analyze\_article      & B     & Article analysis \\
POST   & /api/synthesize            & C     & RAG + LLM chatbot \\
GET    & /api/learning/next-card    & Learn & BKT learning card \\
POST   & /api/learning/submit-ans.  & Learn & Answer + BKT update \\
GET    & /api/learning/roadmap      & Learn & Roadmap data \\
POST   & /api/learning/complete-les.& Learn & Complete lesson \\
GET    & /api/learning/stats        & Learn & User statistics \\
\bottomrule
\end{tabular}
\end{table}

\subsection{Authentication Flow}

The authentication system implements a dual-mode strategy. In production, Firebase Admin SDK verifies JWT tokens by decoding and validating the Firebase ID token from the Authorization header. The system handles \texttt{InvalidIdTokenError} (malformed tokens) and \texttt{ExpiredIdTokenError} (expired sessions) with appropriate 401~responses. Credential loading supports two modes: full JSON content from the \texttt{FIREBASE\_ADMIN\_SDK\_JSON} environment variable for cloud deployments, and file-path-based loading from \texttt{FIREBASE\_ADMIN\_SDK\_PATH} for local development. An optional authentication endpoint (\texttt{get\_optional\_user()}) returns \texttt{None} for unauthenticated requests, enabling public endpoint access.

\subsection{Database Schema}

The persistence layer uses four primary tables. The \texttt{transactions} table stores parsed financial data (user\_id, amount, merchant, category, currency, timestamp, raw\_message) with an index on \texttt{user\_id}. The \texttt{risk\_logs} table records risk assessments with both heuristic and ML classifications plus confidence scores. The \texttt{belief\_states} table maintains Bayesian knowledge states per user-concept pair with a composite unique constraint on \texttt{(user\_id, concept\_id)}. The \texttt{interaction\_events} table logs every quiz attempt with correctness, latency, and timestamps for learning analytics.

The dual-database strategy routes reads and writes to Supabase PostgreSQL when the \texttt{SUPABASE\_URL} environment variable is configured, falling back to a local SQLite file with in-memory caching otherwise. This enables seamless transition between development (SQLite) and production (Supabase) without code changes.

% ══════════════════════════════════════════════════════════════════════
%  VII — DESIGN PATTERNS & SE PRINCIPLES
% ══════════════════════════════════════════════════════════════════════
\section{Design Patterns \& SE Principles}

Table~\ref{tab:patterns} summarises the software engineering patterns applied throughout the system.

\begin{table}[!t]
\centering
\caption{Design Patterns Applied}
\label{tab:patterns}
\renewcommand{\arraystretch}{1.15}
\footnotesize
\begin{tabular}{p{2cm} p{5.2cm}}
\toprule
\textbf{Pattern} & \textbf{Application} \\
\midrule
Multi-Agent   & Three autonomous agents with strict separation of concerns \\
Two-Tier      & ML (Tier~1) + LLM (Tier~2) with automatic fallback \\
Singleton     & GroqClient, ConceptGraph, FinBERT pipeline, Supabase client \\
Strategy      & Dual-DB routing (SQLite vs Supabase) via env variable \\
Observer      & NotificationListenerService monitors notification stream \\
DAG           & Concept prerequisite graph with DFS cycle detection \\
Bayesian      & Probabilistic knowledge state (not deterministic scores) \\
RAG           & TF-IDF retrieval + tag-bonus grounds all LLM responses \\
\bottomrule
\end{tabular}
\end{table}

\subsection{Graceful Degradation Architecture}

Graceful degradation is a first-class architectural principle, not an afterthought. Every feature depending on an external API has a deterministic fallback path, organised across three tiers:

\textbf{Tier~1 --- Deterministic Fallbacks:} Regex-based SMS parsing, rule-based risk explanations, keyword-based sentiment analysis with 70+ finance domain keywords, and hardcoded FAQ responses for 6~core financial topics.

\textbf{Tier~2 --- ML Model Fallbacks:} When FinBERT is unavailable, keyword-based sentiment with strong-phrase detection (3$\times$ priority) provides finance-aware scoring. When category models are missing, a 40+ entry keyword-to-category mapping handles classification.

\textbf{Tier~3 --- Cached Fallbacks:} Previous API responses are cached for reuse. Learning cards are cached per mastery level. NSE market data is pre-warmed at startup and served from cache during network failures. News feeds fall back to a 10-article demo feed.

The system never returns an empty or error state to the user. All LLM calls are wrapped in try/except blocks returning \texttt{None} on failure, with callers implementing deterministic alternatives.

\subsection{Performance Optimisation}

Several techniques ensure responsive user experience:
\begin{itemize}[leftmargin=*]
  \item \textbf{Background Model Loading:} FinBERT model is preloaded in a daemon thread at startup, preventing first-request latency.
  \item \textbf{Async I/O:} All FastAPI endpoints use \texttt{async/await} patterns with \texttt{httpx.AsyncClient} for non-blocking LLM and external API calls.
  \item \textbf{In-Memory Caching:} Belief states, learning cards, and served card references are cached in memory for sub-millisecond access, with database persistence for durability.
  \item \textbf{Lifespan Context Manager:} FastAPI 0.93+ lifespan events handle graceful initialisation and shutdown of background resources.
\end{itemize}

% ══════════════════════════════════════════════════════════════════════
%  VIII — EVALUATION & PERFORMANCE
% ══════════════════════════════════════════════════════════════════════
\section{Evaluation \& Performance}

\subsection{System Performance}

Table~\ref{tab:perf} reports observed latencies for critical operations across the system pipeline.

\begin{table}[!t]
\centering
\caption{System Performance Metrics}
\label{tab:perf}
\renewcommand{\arraystretch}{1.15}
\footnotesize
\begin{tabular}{p{3.2cm} p{1.8cm} p{2cm}}
\toprule
\textbf{Operation} & \textbf{Latency} & \textbf{Component} \\
\midrule
SMS regex parsing      & $<1$ms           & Agent A \\
Category ML inference  & $<10$ms          & Agent A \\
Risk RF prediction     & $<10$ms          & Agent A \\
RAG TF-IDF retrieval   & $<5$ms           & Agent C \\
FinBERT sentiment      & 2--5s / 100      & Agent B \\
Groq LLM call          & 2--8s            & All Agents \\
BKT belief update      & $<1$ms           & Learning \\
Card cache lookup      & $<1$ms           & Learning \\
\bottomrule
\end{tabular}
\end{table}

The Two-Tier architecture ensures that Tier~1 operations (SMS parsing, ML classification, risk assessment, RAG retrieval, BKT updates) complete within milliseconds, providing fast baseline responses. Tier~2 LLM operations add 2--8s of latency but deliver richer, personalised outputs. When Tier~2 is unavailable, users experience no degradation in response time.

\subsection{Feature Availability}

The graceful degradation architecture achieves 100\% core feature availability. All 17~API endpoints return meaningful responses regardless of external service status. When Groq LLM, FinBERT models, or network connectivity are unavailable, the system falls back through the tiered degradation hierarchy. This was validated by testing each endpoint under three conditions: normal operation, LLM service outage, and complete network isolation. In all scenarios, endpoints returned valid, actionable responses with degraded but functional output.

\subsection{Coverage Analysis}

The system provides comprehensive coverage across the personal finance lifecycle:
\begin{itemize}[leftmargin=*]
  \item \textbf{Data Ingestion:} Supports major Indian bank SMS formats (SBI, HDFC, ICICI, Axis) and UPI notification patterns across 10~spending categories.
  \item \textbf{Risk Assessment:} 6-dimensional feature vector captures income, fixed/variable expenses, savings, and derived ratios with month-proration for temporal robustness.
  \item \textbf{Market Intelligence:} Real-time sentiment across 7~market sectors with entity-level granularity and directional trend prediction.
  \item \textbf{Knowledge Coverage:} 32~RAG chunks across 12~financial topic areas spanning beginner to advanced concepts.
  \item \textbf{Learning Path:} 11~concepts across 4~difficulty levels with 5~structured lessons totalling 36~learning cards.
\end{itemize}

% ══════════════════════════════════════════════════════════════════════
%  IX — NOVELTY & KEY CONTRIBUTIONS
% ══════════════════════════════════════════════════════════════════════
\section{Novelty \& Key Contributions}

\subsection{Two-Tier Agentic Pipeline}
Unlike conventional systems that rely solely on ML models or LLMs, every agent operates on a Two-Tier Pipeline. Tier~1 provides fast, deterministic ML/rule-based inference that is always available. Tier~2 layers LLM-powered agentic reasoning on top. If the LLM fails, Tier~1 automatically takes over, guaranteeing \textbf{100\% core feature availability} regardless of external API status---a critical requirement for financial applications where inconsistent service can erode user trust.

\subsection{Unified Personal + Market Context}
No existing consumer product unifies a user's personal transaction data, real-time news sentiment, and curated financial knowledge into a single agentic response. Agent~C injects Agent~A's financial profile (income, expenses, savings rate, risk level), Agent~B's sentiment signals (sector-level bullish/bearish indicators), and RAG-retrieved domain knowledge into every LLM call, producing advice that references the user's actual amounts and current market conditions simultaneously.

\subsection{Bayesian Adaptive Learning for Finance}
The micro-learning engine applies Bayesian Knowledge Tracing---proven in educational technology---to personal finance education. The 3-state belief model with latency-aware updates provides more nuanced learner assessment than binary correct/incorrect scoring. The curriculum compiler dynamically prioritises concepts based on readiness, urgency, and relevance to the user's financial profile, creating a personalised learning path that adapts to both knowledge gaps and real-world financial needs.

\subsection{Backend-Authoritative SMS Parsing}
The Android client sends \emph{all} notifications to the backend without client-side filtering, placing parsing intelligence entirely server-side. This enables rapid iteration on parsing rules without app store updates, ensures a single source of truth for classification logic, and provides consistent behaviour across all device manufacturers---a significant advantage given Android's fragmented notification handling.

\subsection{Domain-Specific NLP Stack}
Rather than a general-purpose sentiment model, the system employs \textbf{FinBERT}---a BERT model fine-tuned on financial text. Combined with 13~boilerplate-removal patterns, 40+ positive and 30+ negative domain keyword fallbacks, strong-phrase detection with 3$\times$ priority weighting, and spaCy NER with a 25+ entry company-to-sector mapping, it achieves finance-aware sentiment accuracy that generic models cannot match.

\subsection{Profile-Aware Curriculum Relevance}
The curriculum compiler's relevance multiplier links learning priorities to the user's actual financial profile. A user with high spending risk receives elevated priority for budgeting and emergency fund concepts ($1.5\times$); a user with debt sees debt management concepts prioritised ($1.6\times$). This profile-financial-education coupling ensures learning has immediate practical applicability.

% ══════════════════════════════════════════════════════════════════════
%  X — FUTURE SCOPE
% ══════════════════════════════════════════════════════════════════════
\section{Future Scope}

\subsection{AI \& Model Enhancements}
Future work includes replacing custom TF-IDF RAG with FAISS or Pinecone vector stores for semantic similarity retrieval, adding a self-evaluation loop where Agent~C critiques its own output before delivery, incorporating reinforcement learning from user feedback to improve recommendation quality, and introducing LSTM or Prophet time-series models for predictive spending forecasting based on historical transaction patterns.

\subsection{Product Features}
Planned features include voice-based financial assistant via speech-to-text integration, portfolio tracking with Zerodha or Groww brokerage APIs for investment monitoring, multi-language SMS parsing for Hindi and regional Indian languages to expand accessibility, and Play Store release with cross-platform iOS support using Kotlin Multiplatform.

\subsection{Engineering \& Operations}
Engineering improvements target MLflow experiment tracking and model registry for ML lifecycle management, a formal evaluation framework incorporating accuracy, AUC-ROC, and human judgement metrics, CI/CD pipeline with GitHub Actions for automated testing and deployment, and on-device SMS parsing using TensorFlow Lite for enhanced privacy by keeping financial data on the user's device.

% ══════════════════════════════════════════════════════════════════════
%  XI — CONCLUSION
% ══════════════════════════════════════════════════════════════════════
\section{Conclusion}

The Agentic Finance System demonstrates that a multi-agent, two-tier architecture can deliver personalised, grounded financial guidance by unifying personal transaction data, real-time market sentiment, and curated domain knowledge. With 3~AI agents, 17~API endpoints, 6~ML/AI models, a 32-chunk RAG corpus, Bayesian adaptive learning across 11~financial concepts with 4~difficulty levels, and a full Android client with 30~Kotlin source files, the system covers the end-to-end pipeline from raw SMS ingestion to explainable financial advice.

The deliberate emphasis on graceful degradation---with three tiers of deterministic, ML-based, and cached fallbacks---ensures that every feature remains operational even when external services are unavailable. The profile-aware curriculum compiler bridges the gap between financial data and financial education, ensuring users learn concepts most relevant to their actual financial situation.

The system validates that multi-agent orchestration with deterministic fallbacks can deliver production-grade financial intelligence accessible to students and early-career earners---a demographic traditionally underserved by both financial advisors and existing fintech applications.

% ── References ────────────────────────────────────────────────────────
\begin{thebibliography}{10}

\bibitem{finbert}
D.~Araci, ``FinBERT: Financial sentiment analysis with pre-trained language models,'' \textit{arXiv preprint arXiv:1908.10063}, 2019.

\bibitem{bkt}
A.~T.~Corbett and J.~T.~Anderson, ``Knowledge tracing: Modeling the acquisition of procedural knowledge,'' \textit{User Modeling and User-Adapted Interaction}, vol.~4, no.~4, pp.~253--278, 1994.

\bibitem{rag}
P.~Lewis \textit{et al.}, ``Retrieval-augmented generation for knowledge-intensive NLP tasks,'' in \textit{Proc.\ NeurIPS}, 2020.

\bibitem{fastapi}
S.~Ram\'{i}rez, ``FastAPI: Modern, fast web framework for building APIs with Python 3.6+,'' 2019. [Online]. Available: https://fastapi.tiangolo.com

\bibitem{llama}
H.~Touvron \textit{et al.}, ``LLaMA: Open and efficient foundation language models,'' \textit{arXiv preprint arXiv:2302.13971}, 2023.

\bibitem{jetpack}
Google, ``Jetpack Compose: Android's modern toolkit for building native UI,'' 2021. [Online]. Available: https://developer.android.com/jetpack/compose

\bibitem{spacy}
M.~Honnibal and I.~Montani, ``spaCy 2: Natural language understanding with Bloom embeddings, convolutional neural networks and incremental parsing,'' 2017.

\bibitem{sklearn}
F.~Pedregosa \textit{et al.}, ``Scikit-learn: Machine learning in Python,'' \textit{Journal of Machine Learning Research}, vol.~12, pp.~2825--2830, 2011.

\bibitem{multiagent}
M.~Wooldridge, \textit{An Introduction to MultiAgent Systems}, 2nd~ed.\hspace{0.6em} Wiley, 2009.

\bibitem{upi}
NPCI, ``Unified Payments Interface: Product overview,'' National Payments Corporation of India, 2023. [Online]. Available: https://www.npci.org.in/what-we-do/upi/product-overview

\end{thebibliography}

\end{document}
